
\subsection{电磁应用}

\begin{tabularx}{\columnwidth}{XX}
  速度选择器&$qE=qvB$\\
  选择速度&$v=\dfrac{E}{B}$\\
  加速电场&$qU=\dfrac12mv^2$\\
  质谱仪半径&$R=\dfrac1B\sqrt{\dfrac{2mU}{q}}$\\
  霍尔电压&$U_H=\dfrac1{nq}\dfrac{IB}{d}$\\
  回旋加速器最大速度&$v_{\max}=\dfrac{qBR}{m}$\\
  回旋加速器最大动能&$E_{k\max}=\dfrac{q^2B^2R^2}{2m}$\\
\end{tabularx}

\begin{formulanotebox}
速度选择器只能选择特定方向进入的粒子，不能单独分辨电性；电性改变时电场力和洛伦兹力方向会同时反向。

回旋加速器的最大动能由磁场和 D 形盒半径决定，与加速电压无关；加速电压主要影响加速次数和总时间。
\end{formulanotebox}

\begin{keypointbox}[title={\strut\textcolor{white}{\bfseries 重点知识点：电磁应用}}]
\textbf{质谱仪}：粒子先经电压 $U$ 加速，再进入匀强磁场偏转。半径越大，$\frac{m}{q}$ 越大；比荷 $\frac{q}{m}$ 越大，半径越小。\\

\textbf{霍尔效应}：载流子在磁场中受洛伦兹力偏转，积累电荷形成横向电压。霍尔电压正负可用于判断载流子电性。\\

\textbf{回旋加速器}：粒子每半周被电场加速一次，电场频率要与粒子回旋频率匹配。
\end{keypointbox}

\begin{casetypebox}[title={\strut\textcolor{white}{\bfseries 常见题型：电磁组合装置}}]
\begin{itemize}[leftmargin=1.5em]
  \item 速度选择器先用 $v=\frac{E}{B}$，再把速度带入后续磁偏转半径。
  \item 质谱仪常联立 $qU=\frac12mv^2$ 与 $R=\frac{mv}{qB}$。
  \item 霍尔题先判断载流子偏转方向，再确定哪侧电势高。
  \item 回旋加速器中 $T=\frac{2\pi m}{qB}$，电场换向周期应与粒子半周运动同步。
\end{itemize}
\end{casetypebox}
